\documentclass[11pt,letterpaper]{article}
\usepackage[T1]{fontenc}
\usepackage{lmodern,amsmath,amssymb,amsthm,microtype}
\usepackage[margin=1in]{geometry}
\usepackage[hidelinks]{hyperref}
\providecommand{\tightlist}{\setlength{\itemsep}{0pt}\setlength{\parskip}{0pt}}
\DeclareUnicodeCharacter{2264}{\ensuremath{\le}}
\DeclareUnicodeCharacter{2265}{\ensuremath{\ge}}
\DeclareUnicodeCharacter{2260}{\ensuremath{\ne}}
\DeclareUnicodeCharacter{2212}{\ensuremath{-}}

\DeclareUnicodeCharacter{2211}{\ensuremath{\sum}}

\title{Six-point extensions of \(G_{449}\) with a point outside the denominator-seven grid}
\author{Aafi Khan Malek\\ \small Independent researcher, Ahmedabad, India\\
\small\href{mailto:aafimalek2032@gmail.com}{aafimalek2032@gmail.com}}
\date{28 November 2026}
\begin{document}
\maketitle
\textbf{Theorem.} Let \(G_{449}\) be the Euclidean unit-distance graph in
the \href{https://arxiv.org/abs/2611.28518}{published Heule--Przybocki embedding}. Put
\(L=(1/7)\mathbb Z[\zeta_{42}]\), with \(\zeta_{42}=\exp(2\pi i/42)\).
Adding six distinct new points, at least one outside \(L\), preserves
five-colorability. The proof uses
\href{https://github.com/Aafimalek/novel-discoveries/tree/234e0ce194fdfa544d4b551ccea9898b2239dc69/hadwiger-nelson/verification/g449/prime-seven-catalogue/five-point-extensions}{the five-point extension theorem},
new geometric reductions, and exact integer checks. It is not
Lean-verified and does not change the chromatic bounds for the plane.

\section{Geometry and mixed lemmas}\label{geometry-and-mixed-lemmas}

The previous exact checks give (I) every off-grid point has at most two
original neighbors, (II) one with two original neighbors is not adjacent
to any new grid vertex with at least three original neighbors, and (III)
one with at least one original neighbor is not adjacent to two new grid
vertices with at least four original neighbors. For (II) see
\href{https://github.com/Aafimalek/novel-discoveries/blob/1c624e99b4c05f24b89d26ee941a62d7bd959fc0/hadwiger-nelson/verification/g449/prime-seven-catalogue/four-five-point-extensions/classify_offgrid_four.py}{the
four-point checker}; for (III) see
\href{https://github.com/Aafimalek/novel-discoveries/blob/1c624e99b4c05f24b89d26ee941a62d7bd959fc0/hadwiger-nelson/verification/g449/prime-seven-catalogue/four-five-point-extensions/one_base_high_pair.py}{the
mixed-five-point checker}. Any two points have at most two common unit
neighbors. The grid neighbors of an off-grid point are independent: both
equilateral apices of a grid unit segment lie in the grid.

A fourth exact observation is (IV): an off-grid point cannot be adjacent
to three new grid vertices with at least four original neighbors. There
are 1,043 such grid vertices. The grid direction catalogue produces all
grid common centers of their triples. Of 188,560,841 unordered triples,
186,404,358 are pairwise nonadjacent, the only triples that could have
an off-grid center. For a triple on a unit circle the scaled
circumradius identity \(UVW+49D^2=0\) is necessary. The same four
justified residue checks as before (43, 127, 211, 337) leave exactly
133,371 triples, all with a verified grid center. Uniqueness of a
noncollinear circumcenter proves (IV).

\section{One off-grid point among
six}\label{one-off-grid-point-among-six}

In a failing six-point extension, deleting any one new point leaves a
five-colorable graph. Every new vertex thus has total degree at least
five. By (I), an off-grid point \(x\) has at least three new neighbors.
If there is just one off-grid point, its new neighbors are independent.
\begin{itemize}
\item If there are three, \(x\) has exactly two base neighbors. By (II) its
three new neighbors each have base degree at most two. Each must then
see both of the remaining grid vertices to achieve total degree at least
five, contrary to the common neighbor bound.
\item If there are four, they
each have base degree at least three (only \(x\) and the remaining grid
point are possible new neighbors), so \(x\) has exactly one base
neighbor by (II). By (III) at least three of them have base degree
exactly three. These three all require adjacency to the remaining grid
point, which would then have three common neighbors with \(x\).
\item If there are five, all five have base degree at least four, contrary to
(IV).
\end{itemize}

\section{Finite shape and prism
pruning}\label{finite-shape-and-prism-pruning}

The shape checker exhausts the \(2^{15}\) possible labeled new graphs on
six points for each fixed number of off-grid labels. At each vertex the
graph induced by its neighbors has degree at most two and no cycle,
since unit edges on its unit circle follow angles \(\pm60^\circ\) and
here that circle has at most five of the six new points. Use the
common-neighbor bound, the independence of grid neighbors of every
off-grid point, (I)--(IV) and the elementary lower bound \(5-d\) on base
degree for any new vertex with new degree \(d\).

Two further geometric rules:
\begin{itemize}
\item Any realized unit triangular prism on six
distinct points consists of two translates of an equilateral triangle,
matching by translation. For an orientation-preserving correspondence,
center the first triangle at zero and write the second \(t+Ra_i\). If
\(R\ne I\), the condition that \(|t+(R-I)a_i|\) be independent of \(i\)
forces \(t=0\), since \(R-I\) is invertible and the \(a_i\) span
affinely. A unit matching distance then means \(R\) rotates by
\(\pm120^\circ\) and the point sets coincide. For an
orientation-reversing correspondence let the normal to the reflection
axis be vertical. Equality of the three matching distances allows at
most two different normal projections of the \(a_i\); two \(a_i\) must
share a projection. Thus the axis is parallel to their side; set them at
\((\pm1/2,0)\) and the third at \((0,\sqrt3/2)\). The matching distances
give \(|t|=|t-(0,\sqrt3)|=1\), so \(t=(\pm1/2,\sqrt3/2)\), again with
coincidences. Hence only translation remains. If one triangle is
entirely in the grid and one point of the other lies in the grid, the
translation is a grid displacement and both triangles are entirely in
the grid.
\item Distinct exterior apices of one equilateral unit triangle,
of forms \(a+b-c\) and \(a+c-b\) up to relabeling, are at distance two.
Their unique common unit neighbor is their midpoint \(a\); no other new
vertex can meet both.
\end{itemize}

In every surviving prism with three or more off-grid points, each grid
vertex has an off-grid neighbor. Every prism vertex has new degree
three, so its off-grid vertices have two original neighbors by minimum
total degree and (I); all adjacent grid vertices have base degree at
most two by (II). With any proper base coloring, all three or more
remaining colors at each vertex suffice for the prism. To prove
3-list-colorability, trim all lists to exactly three. If all lists
agree, assign distinct colors on the first triangle and cyclically shift
them for the matching second triangle. Otherwise choose adjacent \(u,v\)
with different lists and give \(u\) a color absent from \(v\). The prism
minus \(u\) is connected. Color it greedily in reverse spanning-tree
order rooted at \(v\); each nonroot has its parent still uncolored, and
at the root the already colored \(u\) contributes no forbidden color in
its list. After these steps the shape counts for 1 through 6 off-grid
points, up to type-preserving relabeling, are \(0,4,1,0,0,0\).

\section{Remaining two and three off-grid
shapes}\label{remaining-two-and-three-off-grid-shapes}

For exactly two off-grid points, every remaining graph has the two
off-grid points adjacent and sharing two grid vertices whose new degree
is at most three. These two grid vertices thus have base degree at least
two and are the opposite equilateral apices, at distance \(\sqrt3\).
There are 7,056 new grid vertices with base degree at least two. Out of
their 24,890,040 pairs, exact polynomial arithmetic after necessary
modular filters finds 5,691 at this distance. With scaled grid
numerators \(A,B\), the two common unit centers in the same scaling are
\[
(2A+B+\eta(B-A))/3,\qquad (A+2B-\eta(B-A))/3,
\qquad \eta=\zeta_{42}^7.
\] The formula is the midpoint plus or minus a perpendicular offset of
length \(1/2\) for the unscaled segment of length \(\sqrt3\). The
checker finds both numerators coefficientwise divisible by three in
every one of the 5,691 pairs and checks all required distances exactly.
Hence both common centers are grid points, contrary to the remaining
shapes.

For exactly three off-grid points, the only remaining shape is an
off-grid equilateral unit triangle together with its three exterior grid
apices. The apices have new degree two and therefore base degree at
least three. They form a side-two equilateral triangle, whose side
midpoints are the three off-grid vertices. Each off-grid vertex has new
degree four and must have a base neighbor. The 2,429 new grid vertices
with base degree at least three give 2,948,806 pairs. The exact checker
finds 637 distance-two pairs and 21 side-two equilateral triangles; 14
of these have three off-grid side midpoints, and none has a base
neighbor at every midpoint. This closes the last remaining off-grid
shape.


\section{Exact replay summary}
The four supplied Python scripts regenerate the signed unit directions and
the relevant finite grid catalogues from pinned input data. The three
distance searches use only modular filters followed by exact arithmetic.
\begin{center}
\small
\begin{tabular}{r l l}
\textbf{Off-grid points} & \textbf{Residual geometry} & \textbf{Exclusion}\\\hline
1 & Three high-degree grid neighbors & Exact triple check\\
2 & Two opposite grid apices & Common centers on-grid\\
3 & Three exterior grid apices & Missing base neighbor\\
4--6 & Triangular prism & Three-list coloring
\end{tabular}
\end{center}
Adding six all-grid points remains a separate case.
\end{document}

